\section{Sistemi Lineari e Gauss}
\emph{[Rif. Lezione 1, Lezione 2 e 3]}

\subsection{Concetti Introduttivi (Lezione 1)}

\subsubsection{Premesse Fondamentali}
Salvo esplicito avviso contrario, supponiamo sempre di lavorare in $\mathbb{R}$ (insieme dei numeri reali).

\subsubsection{Equazioni Lineari}
\textbf{Definizione:} Un'equazione lineare è un'equazione di I grado (in cui l'incognita compare al grado 1).

\paragraph{Caso 1: Equazioni in una variabile}
Forma normale: $ax + b = 0$.
\begin{itemize}
    \item Se $a \neq 0 \implies$ \textbf{1 soluzione}: $x = -\frac{b}{a}$.
    \item Se $a = 0$ e $b \neq 0 \implies$ \textbf{Impossibile}.
    \item Se $a = 0$ e $b = 0 \implies$ \textbf{Indeterminata} ($\forall x \in \mathbb{R}$).
\end{itemize}

\paragraph{Caso 2: Equazioni in due variabili}
Esempio: $x + y = 1$.
Soluzioni: $\{ (t, 1-t) \mid t \in \mathbb{R} \}$.
\textbf{Interpretazione Geometrica:} I punti di una retta.

\subsection{Sistemi di Equazioni Lineari}
\textbf{Definizione:} Insieme di equazioni che devono essere soddisfatte \textbf{contemporaneamente}.

\subsubsection{Interpretazione Geometrica (Sistemi $2 \times 2$)}
Cercare l'intersezione tra due rette.
\begin{center}
\begin{tabular}{ccc}
\toprule
\textbf{Configurazione} & \textbf{Soluzioni} & \textbf{Esempio} \\
\midrule
\textbf{Rette Incidenti} & \textbf{1 Soluzione} & $\begin{cases} x+y=1 \\ 2x-y=0 \end{cases}$ \\
\midrule
\textbf{Rette Parallele} & \textbf{0 Soluzioni} & $\begin{cases} x+y=1 \\ x+y=2 \end{cases}$ \\
\midrule
\textbf{Rette Coincidenti} & \textbf{$\infty$ Soluzioni} & $\begin{cases} x+y=1 \\ 2x+2y=2 \end{cases}$ \\
\bottomrule
\end{tabular}
\end{center}

\subsection{Generalizzazione}
Un sistema lineare può essere:
\begin{enumerate}
    \item \textbf{Possibile (Compatibile):}
    \begin{itemize}
        \item 1 Soluzione unica.
        \item $\infty$ Soluzioni.
    \end{itemize}
    \item \textbf{Impossibile:} 0 Soluzioni.
\end{enumerate}

\section{Algoritmo di Eliminazione di Gauss}

\subsection{Sistemi Equivalenti}
Due sistemi si dicono \textbf{EQUIVALENTI} se hanno lo stesso insieme di soluzioni.

\begin{concept}{Sistemi Triangolari}
I sistemi triangolari superiori sono facili da risolvere con la \textbf{Sostituzione a Ritroso}.
\end{concept}

\subsection{Le 3 Mosse di Gauss}
Operazioni lecite sulle righe della matrice completa $(A|b)$:
\begin{enumerate}
    \item[(I)] Scambiare due righe ($R_i \leftrightarrow R_j$).
    \item[(II)] Moltiplicare una riga per uno scalare $\lambda \neq 0$ ($R_i \to \lambda R_i$).
    \item[(III)] Sommare a una riga un multiplo di un'altra ($R_i \to R_i + k R_j$).
\end{enumerate}

\subsection{Esempi Pratici}

\subsubsection{Esempio 1: Sistema $3 \times 3$ Completo}
$$ \begin{cases} x+y+z=6 \\ 2x+y-z=3 \\ -x-y+2z=5 \end{cases} \implies C = \left(\begin{array}{ccc|c} 1 & 1 & 1 & 6 \\ 2 & 1 & -1 & 3 \\ -1 & -1 & 2 & 5 \end{array}\right) $$
\textbf{Passo 1:} Uccidere i termini sotto il primo pivot.
\begin{itemize}
    \item $R_2 \to R_2 - 2R_1$: $(0, -1, -3, -9)$
    \item $R_3 \to R_3 + R_1$: $(0, 0, 3, 11)$
\end{itemize}
$$ \left(\begin{array}{ccc|c}  1 & 1 & 1 & 6 \\  0 & -1 & -3 & -9 \\  0 & 0 & 3 & 11  \end{array}\right) $$
Risolvendo a ritroso: $z = 11/3$, $y = -2$, $x = 13/3$.

\subsubsection{Esempio 2: Impossibile}
$$ \left(\begin{array}{cc|c} 1 & 1 & 3 \\ 0 & -4 & -13 \end{array}\right) $$
Se avessimo ottenuto una riga $\begin{pmatrix} 0 & 0 & | & k \end{pmatrix}$ con $k \neq 0$, il sistema sarebbe impossibile.
Esempio:
$$ \left(\begin{array}{ccc|c} 2 & \dots & \dots & \dots \\ 0 & 0 & 0 & 2 \\ 0 & 0 & 0 & 0 \end{array}\right) \implies \text{IMPOSSIBILE} $$

\subsubsection{Esempio 3: Famiglia di Soluzioni}
$$ \left(\begin{array}{ccc|c} 1 & 1 & 1 & 3 \\ 0 & -1 & 1 & 1 \end{array}\right) $$
Ponendo $z=t$ (libero): $y = t-1$, $x = 4-2t$.
Soluzioni: $(4-2t, t-1, t)$ per ogni $t \in \mathbb{R}$.

\subsection{Riepilogo e Morale}
\begin{itemize}
    \item Riga $0 \dots 0 | b \ (b \neq 0) \implies$ \textbf{0 soluzioni}.
    \item Riga $0=0 \implies$ Ridondante.
    \item Sistema quadrato triangolare con pivot $\neq 0 \implies$ \textbf{1 soluzione}.
\end{itemize}
